% !PRINTED 101
\weberpage{101}
From the composition
\weq{7}{\begin{pmatrix}1, & 0\\0, & n\end{pmatrix}\begin{pmatrix}1, & 0\\0, & m\end{pmatrix} = \begin{pmatrix}1, & 0\\0, & mn\end{pmatrix}, \quad \begin{pmatrix}1, & 0\\0, & n\end{pmatrix}\begin{pmatrix}n, & 0\\0, & 1\end{pmatrix} = \begin{pmatrix}n, & 0\\0, & n\end{pmatrix}}
we can conclude further that \gesperrt{every transformation of
degree $n$ can be composed of such transformations as have
prime degree.} If one decomposes $n = pq$ into two factors
 $p$ and $q$ that are relatively prime to one another, then,
determining the numbers $\beta, \delta$ from
\[
p\delta - q\beta = 1
\]
one obtains the composition
\[
\begin{pmatrix}p, & \beta\\q, & \delta\end{pmatrix}\begin{pmatrix}1, & 0\\0, & n\end{pmatrix}\begin{pmatrix}p\delta, & -q\beta\\-1, & 1\end{pmatrix} = \begin{pmatrix}p, & 0\\0, & q\end{pmatrix},
\]
from which it is to be seen that in place of the principal transformation
one may also use any of these transformations $\begin{pmatrix}p, & 0\\0, & q\end{pmatrix}$ for the derivation of all
the others.

\begin{center}
\textbf{§ 30. \gesperrt{The linear fundamental transformations.}}
\end{center}

The whole group $\mathfrak{L}$ of linear transformations can be derived
by the repetition of two among them, which we call the \gesperrt{linear
fundamental transformations}.
If
\[
L = \begin{pmatrix}\alpha, & \beta\\\gamma, & \delta\end{pmatrix}, \quad \alpha\delta - \beta\gamma = 1
\]
is an arbitrary linear transformation, then
\weq{1}{\begin{pmatrix}\alpha, & \beta\\\gamma, & \delta\end{pmatrix}\begin{pmatrix}\delta, & -\beta\\-\gamma, & \alpha\end{pmatrix} = \begin{pmatrix}1, & 0\\0, & 1\end{pmatrix},}
and since the identical transformation is the unit in the group $L$, % [sic] German source prints L; the group is otherwise denoted $\mathfrak{L}$.
\[
L^{-1} = \begin{pmatrix}\delta, & -\beta\\-\gamma, & \alpha\end{pmatrix}
\]
is the \gesperrt{transformation reciprocal to $L$}.

By the power $L^{\pm m}$ we denote that which arises from the $m$-fold
repetition of $L$ or of $L^{-1}$, and we now wish to show
that by the powers of the \gesperrt{fundamental transformations}

\weq{2}{A = \begin{pmatrix}1, & 0\\1, & 1\end{pmatrix}, \quad B = \begin{pmatrix}0, & 1\\-1, & 0\end{pmatrix}}
%CONT>>
